\chapter{Electrons in Crystals: Plane Waves, k-Points, and Metallic Occupations}
\label{ch:solids}

\section{Direct and reciprocal lattices}

A crystal is defined by three primitive vectors $\mathbf a_1, \mathbf a_2,
\mathbf a_3$; the set of translations $\RR = n_1\mathbf a_1 + n_2\mathbf a_2
+ n_3 \mathbf a_3$ ($n_i \in \mathbb Z$) is the Bravais lattice. For FCC gold
the primitive vectors were given in Chapter~\ref{ch:intro}; the primitive
cell volume is $\Omega = |\mathbf a_1 \cdot (\mathbf a_2 \times \mathbf
a_3)| = a^3/4$.

The \textbf{reciprocal lattice} consists of all vectors $\GG$ such that
$\ee^{\ii \GG\cdot\RR} = 1$ for every lattice translation $\RR$. Its
primitive vectors are
\begin{equation}
\mathbf b_1 = \frac{2\pi\, (\mathbf a_2 \times \mathbf a_3)}{\Omega},
\qquad
\mathbf b_2 = \frac{2\pi\, (\mathbf a_3 \times \mathbf a_1)}{\Omega},
\qquad
\mathbf b_3 = \frac{2\pi\, (\mathbf a_1 \times \mathbf a_2)}{\Omega},
\qquad
\mathbf a_i \cdot \mathbf b_j = 2\pi\, \delta_{ij}.
\end{equation}
(In matrix language: if $A$ has rows $\mathbf a_i$ then $B = 2\pi
(A^{-1})^{\mathsf T}$ has rows $\mathbf b_i$ --- implement it that way.) The
reciprocal of FCC is BCC. The Wigner--Seitz cell of the reciprocal lattice is
the \textbf{first Brillouin zone} (BZ); for FCC it is the truncated
octahedron with the high-symmetry points $\Gamma = (0,0,0)$,
$X = \tfrac{2\pi}{a}(1,0,0)$, $L = \tfrac{\pi}{a}(1,1,1)$,
$W = \tfrac{2\pi}{a}(1,\tfrac12,0)$, $K = \tfrac{2\pi}{a}(\tfrac34,
\tfrac34, 0)$, along which you will plot gold's band structure.

\section{Bloch's theorem}

The Kohn--Sham potential of a crystal is lattice-periodic,
$\veff(\rr + \RR) = \veff(\rr)$.

\begin{theorem}[Bloch]
The eigenstates of a Hamiltonian with a lattice-periodic potential can be
chosen as
\begin{equation}
\psi_{n\kk}(\rr) = \ee^{\ii \kk\cdot\rr}\, u_{n\kk}(\rr),
\qquad
u_{n\kk}(\rr + \RR) = u_{n\kk}(\rr),
\label{eq:bloch}
\end{equation}
labeled by a wavevector $\kk$ in the first Brillouin zone and a discrete band
index $n$.
\end{theorem}

\begin{proof}
Let $\hat T_{\RR}$ be the translation operator, $(\hat T_{\RR}\psi)(\rr) =
\psi(\rr + \RR)$. Periodicity of the potential means $[\Ha, \hat T_{\RR}] =
0$ for all $\RR$, and the $\hat T_{\RR}$ commute among themselves, so all can
be simultaneously diagonalized with $\Ha$. Since $\hat T_{\RR}$ is unitary
its eigenvalues have modulus one, and the group law $\hat T_{\RR}\hat
T_{\RR'} = \hat T_{\RR + \RR'}$ forces the eigenvalue to be an exponential of
a linear function of $\RR$: $\hat T_{\RR} \psi = \ee^{\ii\kk\cdot\RR}\psi$
for some real vector $\kk$. Defining $u(\rr) \equiv
\ee^{-\ii\kk\cdot\rr}\psi(\rr)$, one checks $u(\rr+\RR) = u(\rr)$. Finally,
$\kk$ and $\kk + \GG$ give the same eigenvalue of every $\hat T_{\RR}$, so
$\kk$ may be restricted to the first BZ.
\end{proof}

Physically: in a perfect crystal, momentum is conserved modulo reciprocal
lattice vectors; $\kk$ (``crystal momentum'') is the corresponding quantum
number. The spectrum at each $\kk$ is discrete ($\varepsilon_{n\kk}$), and
plotted along BZ paths gives the \emph{band structure}. Bloch's theorem is
also a factorization of the computation: \textbf{the giant eigenproblem of
the infinite crystal decouples into independent moderate-size eigenproblems,
one per $\kk$-point} --- the single most important source of parallelism in
your program.

\section{The plane-wave basis}

Since $u_{n\kk}$ is cell-periodic, expand it in the discrete Fourier series
of the lattice:
\begin{equation}
\psi_{n\kk}(\rr) = \frac{1}{\sqrt{\Omega}} \sum_{\GG}
c_{n\kk}(\GG)\; \ee^{\ii(\kk + \GG)\cdot \rr}.
\label{eq:pwexpansion}
\end{equation}
The basis functions $\ket{\kk+\GG}$ are orthonormal over the cell. The
expansion is truncated by keeping every $\GG$ with kinetic energy below a
\textbf{cutoff}:
\begin{equation}
\tfrac12 |\kk + \GG|^2 \leq E_{\mathrm{cut}} .
\label{eq:ecut}
\end{equation}
This is the \emph{only} convergence parameter of the basis: increasing
$E_{\mathrm{cut}}$ enlarges the variational space, so the total energy
decreases monotonically toward its converged value (other quantities ---
eigenvalue differences, lattice constants --- also converge, though not
necessarily monotonically).
The number of plane waves is, by counting $\GG$-points inside a sphere,
\begin{equation}
N_{\mathrm{pw}} \approx \frac{\Omega}{6\pi^2}\,
\big(2E_{\mathrm{cut}}\big)^{3/2}.
\end{equation}
For gold with an HGH pseudopotential, $E_{\mathrm{cut}} \approx
40$--$60$\,Ha; with $\Omega \approx 114$\,Bohr$^3$ this gives $N_{\mathrm{pw}}
\sim 1400$--$2500$ coefficients per band per spin component --- small dense
problems, ideal for GPU batching.

Plane waves have three cardinal virtues: (i) they are unbiased --- no atom-
centered basis-set superposition errors; (ii) matrix elements of the kinetic
operator and of any multiplicative potential are trivial (diagonal, and
convolution, respectively); (iii) the map between the coefficient
representation $c(\GG)$ and the real-space representation $u(\rr)$ on a
uniform grid is the FFT, costing $O(N \log N)$. Their vice --- inability to
resolve the violent oscillations of wavefunctions near nuclei --- is cured by
pseudopotentials (Chapter~\ref{ch:pseudo}).

In this basis the Kohn--Sham equation \eqref{eq:kseq} becomes an ordinary
matrix eigenproblem for the coefficient vector at each $\kk$,
\begin{equation}
\sum_{\GG'} H_{\GG\GG'}(\kk)\, c_{n\kk}(\GG') = \varepsilon_{n\kk}\,
c_{n\kk}(\GG),
\qquad
H_{\GG\GG'}(\kk) = \tfrac12|\kk+\GG|^2\,\delta_{\GG\GG'}
+ V(\GG - \GG') + \dots,
\label{eq:pwmatrix}
\end{equation}
where $V(\GG)$ are the Fourier components of the local potential
(convention \eqref{eq:fourierconv} below) and the dots stand for the
nonlocal pseudopotential of Chapter~\ref{ch:pseudo}. The matrix is never
stored; Chapter~\ref{ch:hamiltonian} shows how to apply it.

\paragraph{Fourier conventions (fix these once, use them everywhere).}
Every cell-periodic function $g(\rr)$ --- density, potential, $u_{n\kk}$ ---
is expanded as
\begin{equation}
g(\rr) = \sum_{\GG} g(\GG)\, \ee^{\ii\GG\cdot\rr},
\qquad
g(\GG) = \frac{1}{\Omega}\int_{\Omega} g(\rr)\, \ee^{-\ii\GG\cdot\rr}\,\dd\rr .
\label{eq:fourierconv}
\end{equation}
With this convention $g(\GG{=}0)$ is the cell average of $g$, and
$\int_\Omega g^* h\,\dd\rr = \Omega\sum_{\GG} g^*(\GG) h(\GG)$. On an
$N_1{\times}N_2{\times}N_3$ grid with points
$\rr_{j} = \sum_i (j_i/N_i)\,\mathbf a_i$, Eq.~\eqref{eq:fourierconv} is
exactly the discrete Fourier transform:
$g(\GG) = \tfrac{1}{N}\,\mathrm{FFT}[g(\rr_j)]$ and
$g(\rr_j) = N\cdot\mathrm{IFFT}[g(\GG)]$ with $N = N_1N_2N_3$ and the usual
library normalization (forward unnormalized, inverse divided by $N$).
Integer FFT frequency $(m_1,m_2,m_3)$ corresponds to
$\GG = \sum_i m_i \mathbf b_i$, with $m_i$ running over
$-N_i/2 \dots N_i/2 - 1$ in the library's wrapped order. For wavefunctions
the coefficients in Eq.~\eqref{eq:pwexpansion} carry an extra
$1/\sqrt\Omega$ so that normalization is simply
$\sum_{\GG,s} |c^s_{n\kk}(\GG)|^2 = 1$; then
$u_{n\kk}(\rr_j) = \tfrac{N}{\sqrt\Omega}\,\mathrm{IFFT}[c]$ and the density
contribution of the state is $|u_{n\kk}(\rr_j)|^2$ in electrons per
Bohr$^3$.

\paragraph{Densities need a doubled sphere.} The density $n(\rr) \sim
|\psi|^2$ contains products of plane waves, hence Fourier components up to
$2\sqrt{2E_{\mathrm{cut}}}$: represent $n$, $\veff$, and all potentials on a
grid/sphere with cutoff $4E_{\mathrm{cut}}$ (``dense grid''). Choosing the
FFT grid at least this large avoids aliasing errors when forming
$n$ and when multiplying $\veff \times \psi$ in real space.

\section{A first band gap: the two-plane-wave model}
\label{sec:nfe}

Before any machinery, solve Eq.~\eqref{eq:pwmatrix} by hand in the one
situation where two plane waves suffice. It explains where band gaps come
from, it is the first nontrivial test of your code (Chapter~\ref{ch:verify},
Stage~1), and --- with one extra observation about symmetry --- it is the
seed of the topological property computed in Chapter~\ref{ch:topology}.

Take a weak local potential and a $\kk$ on the zone boundary, halfway to a
reciprocal lattice vector: $\kk = \tfrac12\GG_1$ (in FCC, $X$ and $L$ are
such points). The two plane waves $\ket{\kk}$ and $\ket{\kk - \GG_1}$ have
wavevectors $+\tfrac12\GG_1$ and $-\tfrac12\GG_1$: the same kinetic energy
$T = \tfrac18|\GG_1|^2$. Degenerate states are mixed by any perturbation
that connects them, however weak, and the potential does so through its
Fourier component $V_1 \equiv V(\GG_1)$. Keeping just these two basis
functions, Eq.~\eqref{eq:pwmatrix} is
\begin{equation}
\begin{pmatrix} T & V_1 \\ V_1^* & T \end{pmatrix}
\begin{pmatrix} c(\mathbf 0) \\ c(-\GG_1)\end{pmatrix}
= \varepsilon \begin{pmatrix} c(\mathbf 0) \\ c(-\GG_1)\end{pmatrix}
\qquad\Longrightarrow\qquad
\varepsilon_\pm = T \pm |V_1| .
\end{equation}
A gap $2|V_1|$ opens: this is the origin of every band gap at a zone
boundary. Now the observation. If the crystal has an atom at the origin and
inversion symmetry (as gold does), $V(\rr) = V(-\rr)$ makes $V_1$ real, and
the eigenvectors are $(1, \pm1)/\sqrt2$ --- in real space
\begin{equation}
\begin{aligned}
\psi_{\mathrm{even}} &\propto \ee^{\ii\GG_1\cdot\rr/2} + \ee^{-\ii\GG_1\cdot\rr/2}
\propto \cos\!\big(\tfrac12\GG_1\cdot\rr\big),
&\qquad \varepsilon_{\mathrm{even}} &= T + V_1,\\
\psi_{\mathrm{odd}} &\propto \ee^{\ii\GG_1\cdot\rr/2} - \ee^{-\ii\GG_1\cdot\rr/2}
\propto \sin\!\big(\tfrac12\GG_1\cdot\rr\big),
&\qquad \varepsilon_{\mathrm{odd}} &= T - V_1 .
\end{aligned}
\label{eq:nfeparity}
\end{equation}
The cosine state is \emph{even} under $\rr \to -\rr$ and piles its charge
on the atoms, like an $s$ orbital; the sine state is \emph{odd}, has a node
on every atom and its charge between them, like a $p$ orbital. Which one
is lower is decided by the sign of a single number:
\begin{itemize}[leftmargin=2em]
\item $V_1 < 0$ (the electron is attracted to the atomic site at this length
scale): the even, $s$-like state is below the odd, $p$-like one --- the
order you would guess from atomic physics, where $s$ levels lie below $p$
levels.
\item $V_1 > 0$: the order is \emph{inverted}; the odd state is lower.
\end{itemize}
Both signs occur in real materials, because the effective potential felt by
valence electrons is not simply attractive (Chapter~\ref{ch:pseudo}), and
because other bands of the same parity can push one of the two states.
Remember this picture: an integer-valued property of gold's band structure
will turn out to be nothing more than a record of which of these two
orderings is realized at $X$ and at $L$.

\section{Brillouin-zone integration}
\label{sec:bzintegration}

Every observable of the crystal is a BZ average; e.g.\ the density and
band energy are
\begin{equation}
n(\rr) = \sum_n \int_{\mathrm{BZ}} \frac{\dd\kk}{\Omega_{\mathrm{BZ}}}\,
f_{n\kk}\, |\psi_{n\kk}(\rr)|^2,
\qquad
E_{\mathrm{band}} = \sum_n \int_{\mathrm{BZ}}
\frac{\dd\kk}{\Omega_{\mathrm{BZ}}}\, f_{n\kk}\, \varepsilon_{n\kk}.
\end{equation}
Numerically the integral becomes a weighted sum over a finite mesh. The
standard \textbf{Monkhorst--Pack} mesh is the uniform grid
\begin{equation}
\kk_{(m_1 m_2 m_3)} = \sum_{i=1}^{3} \frac{2 m_i - N_i - 1}{2 N_i}\,
\mathbf b_i,
\qquad m_i = 1, \dots, N_i,
\label{eq:mpmesh}
\end{equation}
i.e.\ a uniform $N_1{\times}N_2{\times}N_3$ grid centered on the origin. For
odd $N_i$ it contains $\Gamma$; for even $N_i$ it straddles $\Gamma$
symmetrically without containing it (the fractional coordinates are
$\pm\tfrac{1}{2N}, \pm\tfrac{3}{2N}, \dots$). The alternative
``$\Gamma$-centered'' grid uses fractional coordinates $m_i/N_i$,
$m_i = 0,\dots,N_i{-}1$, and always contains $\Gamma$. Both converge to
the same answer as $N_i\to\infty$, but for the FCC lattice they are
\emph{not} equally good at finite $N_i$: the $\Gamma$-centered grid is
mapped onto itself by all 48 operations of the cubic point group, whereas
the even-$N$ grid of Eq.~\eqref{eq:mpmesh}, built along the primitive
$\mathbf b_i$, is not (its offset
$\tfrac{1}{2N}(\mathbf b_1{+}\mathbf b_2{+}\mathbf b_3)$ is not invariant
modulo the grid spacing). A symmetry-breaking mesh produces a density of
slightly lower symmetry, and the damage is visible: in a $4^3$
calculation for gold the fourfold-degenerate $d$ levels at $\Gamma$ split
by $8$\,meV and the four equivalent $L$ points come out with energies
differing by $22$\,meV; with an $8^3$ mesh the same defects are still
about $1$ and $2$\,meV. With a $\Gamma$-centered mesh they vanish to
solver accuracy. \textbf{Use $\Gamma$-centered meshes for
gold.} They have a second advantage: for even $N_i$ they contain all eight
points $\tfrac12(n_1,n_2,n_3)$ that Chapter~\ref{ch:topology} needs.
Point-group symmetry can fold the mesh to the irreducible wedge with integer
weights; \emph{for a first code, skip symmetry reduction} --- use the full
mesh with equal weights $w_{\kk} = 1/N_{\kk}$. It is simpler, impossible to
get subtly wrong, and on a GPU the extra $\kk$-points are nearly free because
they batch perfectly, and symmetry-related points \emph{must} then come out
with identical eigenvalues, which is a useful test. (One reduction is free
and safe: time reversal maps
$\kk \to -\kk$ with identical eigenvalues and density,
Section~\ref{sec:symmetry}, so you may keep one point of each $\pm\kk$ pair
with doubled weight.) Gold needs a dense mesh ($12^3$ to $16^3$ with
smearing) for well-converged total energies, because the integrand jumps
at the Fermi surface; an $8^3$ mesh is already adequate for band
structures and for debugging. Coordinates:
store each $\kk$ as its three fractional components along
$\mathbf b_{1,2,3}$ and convert to Cartesian with $\kk_{\mathrm{cart}} =
\kk_{\mathrm{frac}} B$ (rows of $B$ are the $\mathbf b_i$). The
high-symmetry points quoted above are, in fractional form,
$\Gamma = (0,0,0)$, $X = (0, \tfrac12, \tfrac12)$,
$L = (\tfrac12,\tfrac12,\tfrac12)$,
$W = (\tfrac14,\tfrac12,\tfrac34)$,
$K = (\tfrac38,\tfrac38,\tfrac34)$ for the $\mathbf a_i$ of
Chapter~\ref{ch:intro}.

\section{Metals: fractional occupations and smearing}
\label{sec:smearing}

In an insulator, bands are either full or empty and $f_{n\kk} \in \{0,1\}$
(per spinor state; no factor 2 once we use spinors). In a metal the
occupation at $T=0$ is a step function $f = \theta(\mu -
\varepsilon_{n\kk})$: the integrand of every BZ integral is discontinuous
across the Fermi surface, and a discrete $\kk$-sum converges miserably
($\sim 1/N_{\kk}$, with wild oscillations). Worse, states at the Fermi level
flip between occupied and empty from one SCF iteration to the next, making
self-consistency unstable.

The cure is \textbf{smearing}: replace the step by a smooth function of width
$\sigma$,
\begin{equation}
f_{n\kk} = f\!\left(\frac{\varepsilon_{n\kk} - \mu}{\sigma}\right).
\end{equation}
Two standard choices to implement:
\begin{enumerate}[leftmargin=2em]
\item \textbf{Fermi--Dirac}: $f(x) = 1/(\ee^{x} + 1)$. Physically a finite
electronic temperature $T = \sigma/k_B$. The associated entropy per state is
$s = -[f \ln f + (1-f)\ln(1-f)]$.
\item \textbf{Gaussian / Methfessel--Paxton}: $f_0(x) = \tfrac12
\mathrm{erfc}(x)$, and the MP family adds Hermite-polynomial corrections,
$f_1(x) = f_0(x) - \tfrac{x}{2\sqrt\pi} \ee^{-x^2}$ (order 1), which makes the
smeared free energy insensitive to $\sigma$ to higher order --- the usual
choice for metals like gold, with $\sigma \approx 0.01$\,Ha.
\end{enumerate}

\paragraph{The free energy.} With smearing, the variational quantity is not
$E$ but the \textbf{Mermin free energy}
\begin{equation}
F = E - \sum_{n\kk} w_{\kk}\, \sigma\, s(f_{n\kk}),
\label{eq:mermin}
\end{equation}
which is what must decrease during SCF and what enters total-energy
comparisons at fixed $\sigma$. Here $s$ is evaluated at $x = (\varepsilon_{n\kk}-\mu)/\sigma$. For the
non-thermal schemes $s$ is a ``generalized entropy'' defined so that $F$ is
still variational:
$s_0(x) = \ee^{-x^2}/(2\sqrt{\pi})$ for Gaussian and
$s_1(x) = (1 - 2x^2)\,\ee^{-x^2}/(4\sqrt{\pi})$ for Methfessel--Paxton
order~1 (which can be negative --- it is not a physical entropy). The best
estimate of the $\sigma\to0$ energy is $E_0 \approx \tfrac12 (E + F)$ for
Fermi--Dirac and Gaussian smearing; for Methfessel--Paxton $F$ itself is
already the estimate, which is the point of the scheme.

\paragraph{Finding the Fermi level.} $\mu$ is fixed by charge conservation:
\begin{equation}
N_e = \sum_{n\kk} w_{\kk}\,
f\!\left(\frac{\varepsilon_{n\kk} - \mu}{\sigma}\right).
\label{eq:fermibisect}
\end{equation}
The right-hand side is monotonically increasing in $\mu$ for Fermi--Dirac
and Gaussian smearing (for MP orders $\ge 1$ the occupation function
overshoots and $N(\mu)$ can decrease slightly across gapped regions of the
spectrum; bisection between the brackets below still converges to a valid
root, but do not assume uniqueness when debugging); solve by bisection
(or Newton with bisection fallback) between $\min \varepsilon - 10\sigma$ and
$\max\varepsilon + 10\sigma$ to $\sim 10^{-12}$. Include enough
\emph{empty bands} in the eigensolver: 11-electron gold fills, on average,
11 spinor bands per $\kk$-point, so solve for $\geq 20$ bands so that
the smeared occupations have room to spill and the eigensolver has a buffer
(Chapter~\ref{ch:algorithms}).

\begin{gpubox}{GPU/JAX notes for this chapter}
The $(\kk, n)$ pair is your fundamental batch axis. Store all coefficients as
one array of shape $[N_k, N_{\mathrm{bands}}, 2, N_G^{\max}]$ (spinor
dimension 2), padding each $\kk$-point's sphere to a common $N_G^{\max}$ with
a boolean mask (JAX demands static shapes; masked padding wastes a few
percent and buys perfect batching). Fermi search
\eqref{eq:fermibisect} is a scalar bisection over a reduction of the whole
$\varepsilon$ array --- do it on device with \texttt{lax.while\_loop}; a
fixed 60--80 iteration \texttt{fori\_loop} bisection is even simpler and
JIT-friendly. Smearing functions must be implemented with numerically safe
forms ($\mathrm{erfc}$, \texttt{log1p}, clipping $x$ to $\pm 40$) to avoid
NaNs poisoning gradients.
\end{gpubox}
